\subsection{Jacobson Radical under Change of Rings}
\setcounter{subsection}{10}

\begin{exr}[title=Corollary 5.2] %NOTE
	Let $R$ be a commutative ring.
	Prove that the polynomial ring $R[x]$ is semiprimitive
	(meaning its Jacobson radical is zero) if and only if $R$ is a
	reduced ring (meaning it has no nilpotents).
\end{exr}


\nq
\begin{exr}[title=Theorem 5.3] %NOTE
	Prove that if $k$ is a field and $R$ a finitely-generated $k$-algebra,
	then $R$ is semiprimitive.
\end{exr}


\nq
\begin{exr}[title=Proposition 5.6] %NOTE
	Let $R$ be a ring, and let $A$ be a subring of $R$.
	Suppose that $A$ is a direct summand of $R$ (as a left $R$-module).
	Prove that $J(A)=J(R)\cap A$.
\end{exr}


\nq
\begin{exr}
	Let $R$ be a commutative ring.
	Prove that the intersection of all prime ideals of $R$ is equal to %TODO
	the set of nilpotents in $R$.
\end{exr}


\nq
\begin{exr}
	Let $R$ be a commutative ring and let
	$p(x)=\sum_{i=0} a_ix^i$ be in $R[x]$.
	Prove that $p(x)$ is a unit in $R[x]$ if and only $a_0$ is a unit in $R$, %TODO
	and $a_1,\ldots,a_n$ are nilpotents in $R$.

	Hint: show that $a_1,\ldots,a_n$ belong to every prime ideal of $R$.
	Thus use the previous exercise to deduce that $a_1,\ldots,a_n$ are nilpotent.
\end{exr}


\nq
\begin{exr}[title=Exercise 5.6] %TODO
	Let $R$ be a ring, let $R[[t]]$ be the ring of formal power series in
	an indeterminate $t$ with coefficients from $R$.
	Note that $tr=rt$ for all $r\in R$.
	Let $(t)$ denote the ideal of $R[[t]]$ generated by $t$.
	This coincides with the set of power series with no constant term:
	\begin{equation*}
		(t) \ := \ \{tf(t) : f(t)\in R[[t]]\} \ = \ \{f(t)\in R[[t]] : f(0) = 0\}
	\end{equation*}
	Prove that the Jacobson radical of $R[[t]]$ is:
	\begin{equation*}
		J(R[[t]]) = J(R) + (t)
	\end{equation*}
\end{exr}


\nq
\begin{exr}[title=secretly Theorem 4.14 part b)]
	Prove that if $R$ is artinian and semiprimitive, then $R$ is semisimple.
\end{exr} \

\begin{pf}
	Suppose $R$ is artinian and semiprimitive.
	Note $R$ must have the following two properties: \vspace{-0.2in}
	\begin{enumerate}[a)]
		\item Every nonzero left ideal $U$ contains a minimal left ideal $B$.
			This follows from the fact that $R$ is artinian.
		\item Every minimal left ideal $B$ is a direct summand of $_{R}R$ with
			a maximal left ideal.
	\end{enumerate}
	\begin{block}
		To see property b), note that since $J(R)=0$, then there necessarily exists
		a maximal left ideal $m$ such that $B\nsubseteq m$.
		Since $B\cap m$ is an ideal of $B$ and $B$ is minimal, then $B\cap m=(0)$.
		Since $m$ is maximal, note that $B+m={_{R}R}$, and therefore
		$_{R}R=B\oplus m$.
	\end{block}
	Now, take any minimal left ideal $B_{1}$ and, by b), write:
	\begin{equation*}
		_{R}R \ = \ B_{1}\oplus U_{1}
	\end{equation*}
	Then $U_{1}\neq0$, so by a), there exists a minimal left ideal $B_{2}\subseteq U_{1}$.
	By b), $B_{2}$ is a direct summand of $_{R}R$ and thus $U_{1}$, so we can write:
	\begin{equation*}
		U_{1} = B_{2}\oplus U_{2} \ \implies \ _{R}R = B_{1}\oplus B_{2}\oplus U_{2}
	\end{equation*}
	Continuing ad infinitum, if $R$ is not semisimple, then we necessarily always
	have that $B_{k+1}\subsetneq U_{k}$.
	In other words, we get an infinite descending chain:
	\begin{equation*}
		U_{1} \ \supsetneq \ U_{2} \ \supsetneq \ U_{3} \ \supsetneq \ \cdots
	\end{equation*}
	This contradicts the fact that $R$ is artinian, so we must have that
	$R$ is semiprimitive as needed.
\end{pf}
